\chapter{Tinjauan Pustaka}\label{cha:litreivew}

\section{Perpustakaan Digital (Digital Library)}
Perpustakaan digital adalah sistem informasi yang mengumpulkan, mengelola, dan menyediakan akses ke koleksi digital berupa jurnal, buku, skripsi, dan karya ilmiah lainnya melalui jaringan internet \cite{arms2001digital}. Kehadiran perpustakaan digital di institusi pendidikan mempermudah sivitas akademika dalam mencari referensi tanpa batasan ruang dan waktu \cite{pendit2008perpustakaan}. Penggunaannya sangat bergantung pada kualitas antarmuka dan kemudahan navigasi agar informasi dapat ditemukan secara cepat, tepat, dan akurat oleh pengguna \cite{saleh2020manajemen}.

\section{Sistem Informasi Akademik}
Sistem informasi akademik merupakan integrasi perangkat keras dan perangkat lunak yang dirancang untuk mendukung proses manajemen dan layanan pendidikan di perguruan tinggi \cite{jogiyanto2005analisis}. Layanan ini berfokus pada penyampaian informasi yang efisien kepada mahasiswa dan dosen. Kualitas platform akademik yang baik tidak hanya dinilai dari kelengkapan fiturnya, tetapi juga dari seberapa baik sistem tersebut memenuhi ekspektasi penggunanya dalam menyelesaikan tugas-tugas pencarian literatur \cite{delone2003delone}.

\section{User Interface (UI)}
\gls{ui} adalah bagian visual dari sistem perangkat lunak yang menjadi titik interaksi langsung antara pengguna dan komputer. Desain \gls{ui} mencakup elemen tata letak, warna, tipografi, serta komponen interaktif seperti tombol, kolom pencarian, dan menu navigasi \cite{tidwell2010designing}. Desain \gls{ui} yang baik harus intuitif, terstruktur, dan konsisten sehingga dapat meminimalkan kesalahan pengguna saat mengoperasikan sistem \cite{galitz2007essential}. Antarmuka yang rumit dan tidak terorganisir seringkali menjadi penyebab utama pengguna meninggalkan sebuah platform digital \cite{setiawan2023implementasi}.

\section{User Experience (UX)}
\gls{ux} mencakup seluruh persepsi, emosi, dan respons pengguna yang muncul sebelum, selama, dan setelah berinteraksi dengan sebuah produk atau sistem \cite{iso2019ergonomics}. Berbeda dengan \gls{ui} yang berfokus pada tampilan visual, \gls{ux} lebih menitikberatkan pada alur operasional, kenyamanan, dan nilai yang dirasakan pengguna saat mencapai tujuannya \cite{hassenzahl2006user}. \gls{ux} yang optimal akan menciptakan pengalaman yang mulus (\textit{seamless}), mengurangi hambatan kognitif (\textit{cognitive load}), dan pada akhirnya meningkatkan motivasi pengguna untuk terus menggunakan sistem \cite{garrett2011elements}.

\section{Desain Interaksi}
Desain interaksi berfokus pada penciptaan dialog yang bermakna antara pengguna dan sistem digital \cite{preece2015interaction}. Prinsip ini memastikan bahwa setiap tindakan pengguna, seperti mengunduh dokumen atau melakukan pencarian referensi, mendapatkan umpan balik (\textit{feedback}) yang jelas dari sistem. Penerapan desain interaksi yang baik dalam perpustakaan digital akan membantu pengguna memahami arsitektur informasi dan hierarki data dengan lebih cepat tanpa memerlukan panduan manual yang rumit \cite{cooper2014about}.

\section{Kepuasan Pengguna (User Satisfaction)}
Kepuasan pengguna adalah respons emosional yang timbul ketika sebuah sistem berhasil memenuhi atau melampaui harapan pengguna dalam menyelesaikan suatu tugas \cite{bhattacherjee2001understanding}. Dalam konteks layanan digital kampus, faktor-faktor seperti kecepatan akses, relevansi hasil pencarian, dan kemudahan pengunduhan sangat memengaruhi tingkat kepuasan \cite{mulyadi2023analisis}. Pengguna yang merasa puas cenderung akan terus menggunakan layanan \gls{digilib} tersebut sebagai sumber referensi utama mereka.

\section{User Experience Questionnaire (UEQ)}
\gls{ueq} adalah instrumen evaluasi terstandarisasi yang digunakan untuk mengukur pengalaman pengguna secara komprehensif, cepat, dan reliabel \cite{schrepp2014applying}. Kuesioner ini mengukur kualitas pragmatis (fungsionalitas) dan kualitas hedonis (emosional) pengguna melalui 26 item pertanyaan. Evaluasi ini terbagi dalam 6 dimensi: \textit{attractiveness} (daya tarik), \textit{perspicuity} (kejelasan), \textit{efficiency} (efisiensi), \textit{dependability} (ketepatan), \textit{stimulation} (stimulasi), dan \textit{novelty} (kebaruan). Skala penilaiannya berkisar dari -3 (sangat negatif) hingga +3 (sangat positif), memungkinkan peneliti untuk melakukan evaluasi kuantitatif (\textit{benchmark}) atas persepsi pengguna terhadap produk digital secara menyeluruh.